\chapter{Methodology}
\label{sec:method}

This section defines the concepts the system is built on and the models it uses; Chapters~\ref{sec:architecture}
to~\ref{sec:production} then describe how they are combined.

\section{Glosses, sign lexicons and poses}
\paragraph{Gloss.} A gloss is the written label of a sign, conventionally in capitals (PRAYER, \ar{الصلاة}). A signed
sentence is written as a gloss sequence in the order of the sign language, which differs from the spoken language:
articles and the copula are dropped, and several words may become one sign. Glosses are the interface between text and
signing: the answer is translated into glosses, and each gloss is looked up in the lexicon.

\paragraph{Sign lexicon.} A sign lexicon is a dictionary from glosses to recorded signs. Each entry holds the gloss, a
sign identifier, the source dataset, the recorded pose, a review status (\emph{approved} or \emph{pending}), and flags
for religious terms and fingerspelling letters. One gloss may have several recordings; one is the canonical sign and
another, by a different signer, is kept as a recognition template. The lexicon defines what the system can sign: a word
is signable if, and only if, it can be matched to an entry.

\paragraph{Pose.} A sign is stored not as video but as a pose sequence $P \in \mathbb{R}^{T\times 50\times 3}$: for each
of $T$ frames (25 per second), 8 upper-body joints and 21 landmarks per hand, extracted with MediaPipe Holistic
\cite{mediapipe,mediapipeholistic}. Raw coordinates $\mathbf{x}_{t,j}$ are normalised per clip,
\begin{equation}
\hat{\mathbf{x}}_{t,j} = \frac{\mathbf{x}_{t,j} - \mathbf{x}_{t,\mathrm{neck}}}{\operatorname{median}_t \lVert
\mathbf{x}_{t,\mathrm{Rsh}} - \mathbf{x}_{t,\mathrm{Lsh}} \rVert},
\end{equation}
so that the neck is the origin and the unit is the signer's shoulder width; signs recorded by different people at
different distances then share one coordinate system and can be joined (Appendix~\ref{app:pose}).

\paragraph{Comparing signs.} Two recordings of a sign differ in speed, so they are compared by dynamic time warping
(DTW) \cite{dtw}, which aligns the frames of two sequences $A$ and $B$ monotonically before summing their distances:
\begin{equation}
D(i,j) = \lVert \mathbf{a}_i - \mathbf{b}_j \rVert + \min\{D(i-1,j),\, D(i,j-1),\, D(i-1,j-1)\}.
\end{equation}
The canonical sign of a gloss is the DTW medoid of its recordings (the recording with the smallest total distance to
the others), and the back-translation baseline recognises a sign by its DTW nearest neighbour among templates.

\section{Lexical matching and morphological analysis}
\emph{Lexical matching} compares a word with the lexicon's glosses as written, after orthographic normalisation.
It fails whenever the text uses a form the dictionary does not list: a plural, a conjugated verb, or a word with an
attached conjunction or preposition. \emph{Morphological analysis} recovers the parts of a word (prefixes, stem,
suffixes) and its \emph{lemma}, the dictionary form shared by all its inflections, so that \ar{الصلوات} (the prayers)
finds the sign \ar{الصلاة} (prayer). How much analysis is needed depends on the language (Table~\ref{tab:morph}).

\begin{table}[ht]
\centering\small
\caption{Lexical normalisation and morphological analysis per language, applied in the same way to the corpus, the
retrieval index, the support check and the lexicon lookup.}
\label{tab:morph}
\begin{tabularx}{\textwidth}{L{1.6cm} L{4.2cm} Y}
\toprule
Language & Morphology & Analysis used \\
\midrule
English & light inflection & lower-casing, accent folding, possessive removal; irregular forms (\emph{went} $\rightarrow$
\emph{go}) and suffix rules (\emph{-ies, -es, -s, -ing, -ed, -ly}) \\
Arabic & root-and-pattern words; attached conjunctions, prepositions, article and pronouns; broken plurals &
orthographic normalisation (alef and hamza forms, taa marbuta, diacritics); clitic stripping (\ar{و، ف، ب، ك، ل، ال},
pronoun suffixes); lemmas from the CAMeL Tools morphological analyser \cite{camel}; regular plurals to singular; number
words to digits \\
Turkish & agglutinative: chains of suffixes on a stem & Turkish-specific casing (\emph{I/\i}, \emph{\.I/i}); the first five
letters as the stem, a simple and robust choice for Turkish retrieval \\
Urdu & separate postpositions; plural and oblique suffixes & folding of Arabic-script letter variants
(\ar{ي/ی}, \ar{ك/ک}, \ar{ه/ہ}); removal of the plural and oblique suffixes \ar{وں، یں، ات، ے} \\
\bottomrule
\end{tabularx}
\end{table}

\paragraph{Arabic in detail.} Arabic is the hardest case. A written word such as \ar{وبالصلوات} is \ar{و} (and) +
\ar{ب} (with) + \ar{ال} (the) + \ar{صلوات} (prayers), and \ar{صلوات} is the plural of \ar{صلاة}. CAMeL Tools proposes,
for each word, all its morphological analyses with part of speech and lemma; we keep the noun and adjective lemmas,
most frequent first. Because Arabic is written without short vowels, one spelling can be several words: \ar{كفر}
(disbelief) and the verb \ar{تكفر} (expiates) share letters, and \ar{شر} (evil) is a prefix of \ar{شرك} (polytheism).
Three guards keep analysis from changing the meaning: a word that can only be a verb matches only its exact form; a
two-letter sign matches only after a prefix was removed or when the analyser gives it as the lemma; and a sign is never
matched inside a longer word that keeps an initial alef (\ar{إسلام} is not \ar{سلام}, peace).

\section{Retrieval-augmented generation}
A large language model writes fluent text, but on its own it answers from what it absorbed in training and can state
plausible errors. \emph{Retrieval-augmented generation} (RAG) \cite{rag} separates knowing from writing: a
\emph{retriever} first selects passages from a trusted collection, and the \emph{generator} (the language model) writes
the answer from those passages, citing them. In Isharati the collection is the Qur'an and the authenticated hadith
(Chapter~\ref{sec:knowledge}), and grounding is enforced after generation by checking every sentence against its
citations. Two complementary retrievers are combined.

\paragraph{Lexical retrieval (BM25).} BM25 \cite{bm25} scores a passage $d$ for a query $q$ by the query terms it
contains, weighting rare terms more and long passages less:
\begin{equation}
\mathrm{BM25}(q,d) = \sum_{t \in q} \mathrm{IDF}(t)\,
\frac{f(t,d)\,(k_1 + 1)}{f(t,d) + k_1\left(1 - b + b\,\frac{|d|}{\mathrm{avgdl}}\right)},
\end{equation}
with $f(t,d)$ the frequency of term $t$ in $d$, $|d|$ the passage length, avgdl the average length, and $k_1 = 1.2$,
$b = 0.75$. Terms are the normalised forms of Table~\ref{tab:morph}. BM25 is exact: it finds names and technical terms,
but misses a passage that says the same thing in other words.

\paragraph{Dense retrieval (embeddings).} An \emph{embedding} model maps a text to a vector such that texts with similar
meaning have nearby vectors; relevance is the cosine similarity $\cos(\mathbf{e}_q, \mathbf{e}_d)$ between the query and
passage vectors. Dense retrieval finds paraphrases and related wording (\ar{الخسوف} also finds \ar{كسوف}, both eclipse
terms) that share no word with the query. We use multilingual E5-large \cite{me5}, initialised from XLM-RoBERTa-large
\cite{xlmr}: 24 layers, 1024-dimensional embeddings, 100 languages, trained first by contrastive learning on more than a
billion weakly paired texts and then fine-tuned on labelled retrieval data. Inputs are prefixed with \code{query:} or
\code{passage:} as the model requires, truncated to its 512-token limit, and vectors are unit-normalised. One model
serves all four languages. Every passage of each corpus is embedded once; a question is embedded at query time.

\paragraph{Fusion.} The two rankings are combined by reciprocal rank fusion \cite{rrf} (Chapter~\ref{sec:knowledge}),
which adds $1/(K + \mathrm{rank})$ over the rankings and so rewards passages that both retrievers place high, without
having to calibrate their scores against each other.

\section{Language models}
Language models are used for five tasks: expanding the question into search keywords, choosing the most direct
passage, writing the cited answer, translating the answer into glosses, and (offline) proposing concept mappings for
uncovered words. All calls use temperature 0. The models are open-weight, so the system can be moved to self-hosted
inference; for the prototype they are called through hosted APIs, in a fallback chain: if a model does not answer within
25\,s, the next is tried (Table~\ref{tab:llms}).

\begin{table}[ht]
\centering\small
\caption{Language models used, in fallback order.}
\label{tab:llms}
\begin{tabularx}{\textwidth}{L{3.3cm} L{2.4cm} Y L{2.1cm}}
\toprule
Model & Parameters & Architecture & Access \\
\midrule
Nemotron 3 Super 120B-A12B \cite{nemotron3super} & 120B total, 12B active per token & hybrid Mamba-2 / attention mixture
of experts, 1M-token context; NVIDIA open model licence & NVIDIA NIM \cite{nvidianim} \\
Nemotron 3 Ultra 550B-A55B \cite{nemotron3ultra} & 550B total, 55B active & same family, larger & NVIDIA NIM \\
gpt-oss-120b \cite{gptoss} & 117B total, 5.1B active & mixture of experts, reasoning model; Apache 2.0 & Groq \\
\bottomrule
\end{tabularx}
\end{table}

Nemotron's step-by-step reasoning mode is switched off, because the outputs must be short, structured answers. A reply
that turns out to be the model's reasoning rather than its answer is treated as a failure and passed to the next model.
For gpt-oss, a reasoning model, the token budget is raised so that its reasoning does not exhaust it before the answer.

\section{Lexicon-constrained glossing}
Translation into glosses is done by the language model, not by rules, because sign order and the choice between
synonyms need context. The model is given the answer, the lexicon's gloss list and examples of sign-language order, and
must use only listed glosses; every gloss it returns is checked against the lexicon, so it cannot invent a sign. Words
with no sign are reported as missing rather than fingerspelled. The Arabic glosser adds the morphological matching above,
phrase signs (\ar{محمد رسول الله}), pronoun mapping and ten word-order examples from ArabSign \cite{arabsign}
(Section~\ref{sec:glossing}).
